\section{Introduction}

When several layers share a learned bank of parameters, their effective
weights can be written as $\Theta=AB$: each row of a router $A$ selects a
convex combination of the rows of a jointly learned basis matrix $B$.
The router is useful for optimization and compression, but its entries are
not uniquely determined by the effective weights. To identify the
\emph{current} factors, we therefore need an additional observation.
Even then, a decision based on those factors may differ from one obtained
by clustering the effective weights.

This ambiguity is known in simplex-structured matrix factorization
(SSMF) theory~\citep{vuthanh2023bounded}. For example, the invertible
transform $M_s=(1-s)I+s\ones\ones^\top/K$, $0<s<1$, preserves the product
under $(A,B)\mapsto(AM_s,M_s^{-1}B)$ while moving each router row toward
uniformity. This increases the entropy of nonuniform rows. It preserves
coordinate ordering and argmax assignments; other feasible, nonuniform
changes of basis can change the argmax partition
(Appendix~\ref{app:secondary-theory-statements}). An effective model alone
therefore need not determine coordinate-based summaries or decisions.

We study an additional observation supplied by a specified optimizer:
the instantaneous response of the effective parameters to an arbitrary
effective gradient. With $A=\operatorname{softmax}_{\rm row}(Z/\tau)$,
fixed temperature, and known positive Euclidean learning rates for $Z$ and
$B$, this is a linear map $\mathcal R_{A,B}:G\mapsto-\dot\Theta$.
We observe this map through controlled interventions at a fixed checkpoint,
using the same optimizer metric throughout. It contains more information
than the response to a single task gradient; observing a passive training
trajectory does not supply the same map.

Under positivity and rank assumptions, the response resolves the static ambiguity.
For strictly positive, full-column-rank routers and full-row-rank bases
with the same product, equal response operators imply a common basis
permutation, and conversely. The proof recovers a Gram matrix and a family
of router covariances from the response, then reduces the remaining
ambiguity to simultaneous diagonalization. On uniformly nondegenerate
domains, the same chain gives an explicit inverse-stability bound.

The complete-operator premise can also be replaced by finite, imperfect
observations. If $D-K\ge L$, we design $K$ unit-norm queries from a noisy
observed product. Every pair of factors consistent with the product and
response error budgets is close modulo permutation, with an explicit bound
on the diameter of the \emph{entire} feasible set. The true basis need not
lie in the estimated product subspace, and candidates need not belong to one
local solution branch. The bound describes the information available on the specified domain.
Finding a feasible solution efficiently and verifying the domain's rank
and interior margins from data remain separate problems.

A structural decision may need less information than the full factors. For equal-capacity hard tying, minimizing full
Euclidean-response error reduces exactly to clustering the router rows;
off-diagonal Gram information suffices. This differs both from clustering
the effective weights and from assigning each row to a fixed router
coordinate. We construct strictly interior, full-rank examples in which these decisions
disagree, including equal products with different response-optimal partitions. A fixed squared-loss
gradient flow can separate coefficient assignment from weight clustering; finite
margins can also protect agreement for a specified interval.

We establish the following results:
\begin{enumerate}
\item A complete Euclidean-response stabilizer characterization, an explicit
inverse on nondegenerate domains, and a noisy-product, finite-query bound
for the entire feasible factor set
(Section~\ref{sec:theory}).
\item An exact balanced response-decision objective and its sufficient Gram
information, explicit conflicts with weight-based decisions, and a
constructive gradient-flow crossing with conditional stability bounds
(Section~\ref{sec:sharing-dynamics}).
\item Recovery and training experiments on pretrained models: a fixed recovery
protocol on nine
pretrained models (160M--8B); a separate 18-run, three-model language-training
study; a selected real-model local SGD reproduction; and controlled
regression/precision checks (Section~\ref{sec:experiments}).
\end{enumerate}

The analysis builds on existing results for response operators,
diagonalization, robust inverses, Gram-based clustering, and trajectory
perturbations. We use these tools to derive the identification guarantees
and decision consequences of the present observation model.
The original nine-checkpoint structural test shows no incremental benefit
from recovery over effective-weight clustering. The subsequent decision study asks whether training can reach a
disagreement. No crossing was detected in its ordinary SGD runs under the
specified schedule, although local SGD from a selected state did cross.
This establishes possibility while leaving its frequency open. Throughout,
identification concerns the current optimizer-dependent factorization;
the historical process that produced it is outside the observation model.
